\documentclass[11pt]{article}
\usepackage[margin=1in]{geometry}
\usepackage{amsmath,amssymb,booktabs,microtype}
\usepackage{xcolor}
\usepackage[colorlinks=true,linkcolor=blue,urlcolor=blue]{hyperref}
\newcommand{\Sym}{\operatorname{Sym}}
\newcommand{\cS}{\mathcal S}
\title{Polar and Dual-Polar Schemes: Permutation Formula\\Proof and Exact Verification for $Sp_4(2)$}
\author{Executable companion note}
\date{17 July 2026}
\begin{document}
\maketitle
\begin{abstract}
We prove the fixed-polar-subspace version of the permutation Molien formula
and test two nontrivial actions of the explicitly constructed matrix group
$Sp_4(2)$: isotropic points and maximal totally isotropic $2$-spaces.  The
latter is a rank-two dual-polar scheme.
\end{abstract}

\section{Cycle and orbit formulas}
Let a finite isometry group $G$ act on a finite family $X$ of totally
isotropic or singular subspaces, and put $V=\mathbb C[X]$.  If $c_\ell(g)$
counts $\ell$-cycles on $X$, the determinant of an $\ell$-cycle block is
$1-t^\ell$, so
\begin{equation}\label{eq:master}
\cS_{G,X}(\mathbf t)=\frac1{|G|}\sum_{g\in G}
\prod_{i,\ell}(1-t_i^\ell)^{-c_\ell(g)}.
\end{equation}
Since the monomial basis of $\Sym^d(V)$ is indexed by $d$-multisets in $X$,
Burnside's lemma independently gives
\begin{equation}\label{eq:orbit}
[m_\tau]\cS_{G,X}=\#\left(G\backslash\prod_j\Sym^{\tau_j}(X)\right).
\end{equation}

Define the geometric fixed-point count
\begin{equation}\label{eq:fixed}
F_r^{\rm pol}(g)=\#\{U\in X:g^rU=U\}.
\end{equation}
The universal relation $F_r=\sum_{\ell\mid r}\ell c_\ell$ and Möbius
inversion yield
\begin{equation}\label{eq:mobius}
c_\ell^{\rm pol}(g)=\frac1\ell\sum_{d\mid\ell}
\mu(\ell/d)F_d^{\rm pol}(g).
\end{equation}
Equations \eqref{eq:master}--\eqref{eq:mobius} prove the proposed formula.
Relative to the Grassmann case, the only additional predicate in
\eqref{eq:fixed} is total isotropy or singularity.

For maximal isotropic subspaces in Witt rank $r$, pair orbits are classified
by $r-\dim(U\cap W)$.  Consequently
\[
[m_{(1)}]=1,\qquad [m_{(2)}]=[m_{(1,1)}]=r+1.
\]

\section{Concrete symplectic construction}
On $\mathbb F_2^4$ use the nondegenerate alternating form
\[
\langle x,y\rangle=x_1y_2+x_2y_1+x_3y_4+x_4y_3.
\]
The code enumerates all $20{,}160$ matrices in $GL_4(2)$ and retains exactly
the 720 satisfying $\langle gx,gy\rangle=\langle x,y\rangle$ on basis
vectors.  It separately enumerates reduced-row-echelon subspaces.  There are
15 one-spaces (all isotropic for an alternating form) and 15 two-spaces on
which the form vanishes identically.  Thus both induced representations have
dimension 15, but their vertices are different geometric objects.

For every group matrix and every needed power, the program acts on the stored
bases and counts invariant members of $X$ directly.  Equation
\eqref{eq:mobius} must reproduce the cycles of the separately constructed
vertex permutation.

\section{Three independent coefficient routes}
For the actual $15\times15$ permutation matrix $P$, complete symmetric
characters are computed without cycle data from
\[
d h_d(P)=\sum_{a=1}^d\operatorname{tr}(P^a)h_{d-a}(P),\qquad h_0=1.
\]
The matrix-group target is the average of $\prod_jh_{\tau_j}(P)$.  A second
calculation extracts the coefficient from \eqref{eq:master}; a third explicitly
enumerates the configuration orbits in \eqref{eq:orbit}.  Agreement of these
routes tests the representation, the proposed formula, and its geometric
interpretation separately.

\section{Exact results}
Values are ordered by
$\tau=(1),(2),(1,1),(3),(2,1),(1,1,1)$.
\begin{center}
\begin{tabular}{@{}lrrl@{}}
\toprule
Action & $|G|$ & $|X|$ & common coefficient vector\\
\midrule
$Sp_4(2)$ on isotropic points & 720 & 15 & $(1,3,3,8,12,16)$\\
$Sp_4(2)$ on maximal isotropic 2-spaces & 720 & 15 & $(1,3,3,8,12,16)$\\
\bottomrule
\end{tabular}
\end{center}
All 12 matrix/cycle/orbit comparisons and all geometric fixed-point
reconstructions pass.  The dual-polar pair value is $3=r+1$ for $r=2$.  For
the maximal action at $(t_1,t_2)=(0.037,0.061)$, both numerical averages equal
$1.1263214575673286$ (reported error $0$).

\section{Scope and reproduction}
The exact proof applies equally to alternating, Hermitian, and quadratic
forms.  The executable representative is symplectic because $Sp_4(2)$ is
large enough to exhibit all rank-two pair relations while remaining fully
enumerable and auditable.

Run
\begin{verbatim}
marimo run marimo_notebooks/polar_scheme.py
\end{verbatim}
from the repository root to inspect every comparison.
\end{document}
